

\song\xiaosi 自动驾驶规划系统面临计算效率、可解释性合规和长尾泛化三大挑战。当前端到端大模型参数量巨大、不可解释，难以满足边缘部署和安全认证需求。神经电路策略（NCP）以极少参数和生物启发的稀疏结构提供了一条差异化路径，但其布线拓扑完全依赖手工设计，且在强化学习场景下的适用性尚未得到验证。

本文提出一种基于NCP的可解释自动驾驶规划方法。首先，将NCP集成为Double DQN的Q-network，针对NCP的ODE约束导致的表达能力瓶颈，提出两层Q-head方案，在保持液态动力学可解释性的同时实现有效的Q值估计。其次，设计面向NCP布线拓扑的进化搜索算法，采用安全感知的多目标适应度函数，自动发现兼顾性能与安全性的最优布线方案。

在Highway-env的4个驾驶场景上，与MLP、LSTM、GRU、全连接LTC等5种基线模型进行了系统对比实验。结果表明，NCP以7,152个参数在高速公路场景上取得了与13,189参数的MLP相当的奖励（34.94对34.88），参数效率是GRU的3倍以上，且方差最低（1.21）。与同参数量的全连接LTC相比，NCP的结构化稀疏拓扑在奖励上提升了23\%。多层次可解释性分析揭示了NCP命令层神经元与驾驶决策之间的对应关系。

\vspace{\baselineskip}
\noindent \xiaosi\hei 关键词: \song\xiaosi 神经电路策略；液态时间常数网络；自动驾驶规划；拓扑搜索；可解释人工智能